\section{The source invariant and its exact polar divisor}\label{sec:integral}

We now use the proved matrix integral of Joshi--Roffelsen. Its algebraic
construction is an external dynamical input; the arguments here identify
its polar divisor on the complete resolved surface and control its
extension at every ordinary point.

\subsection{The explicit integral and its specialization}

For a primitive \(r\)-th root parameter, set
\(R=\ZZ[s]/(\Phi_r(s))\), where \(\Phi_r\) is the cyclotomic
polynomial. At gauge \(w=1\), write
\begin{equation}\label{eq:source-matrix}
A(z)=A_0+zA_1+z^2A_2,
\end{equation}
where
\begin{align*}
A_0&=\begin{pmatrix}
t+x-xy&-x\\
t+x-ty-2xy+xy^2&x(y-1)
\end{pmatrix},\\[3pt]
A_1&=\begin{pmatrix}
y-x+x/y-1-t/x&1\\
y-2x-1+xy+x/y-t/x&1
\end{pmatrix},\qquad
A_2=\begin{pmatrix}1&0\\0&0\end{pmatrix}.
\end{align*}
Define
\begin{equation}\label{eq:source-integral}
I_{r,t}(x,y)=\Tr\bigl(A(s^{r-1})\cdots A(s)A(1)\bigr)-(t^r+1).
\end{equation}
The source proves that this is a Laurent polynomial in \(x,y\),
polynomial in \(t\), with coefficients in \(R\), and that
\begin{equation}\label{eq:invariance}
I_{r,st}\bigl(F_t(x,y)\bigr)=I_{r,t}(x,y)
\end{equation}
as a rational identity over the cyclotomic field
\citep[Theorem 3.1]{JoshiRoffelsen2026}. The trace is independent of
the auxiliary gauge. The source's proof uses the cyclic trace of
\(A(s^{r-1}z)\cdots A(z)\) and telescoping Lax compatibility;
its constant and leading trace coefficients are \(t^r\) and \(1\).
It does not assume the source's later conjecture about genera of fibres.

For the finite field \(k=\F_q\), the order \(r\) divides
\(q-1\), so \(p=\operatorname{char}k\) does not divide \(r\).
An element of exact order \(r\) in \(k^*\) is a zero of
\(\Phi_r\) modulo \(p\). Indeed the factorization
\[
X^r-1=\prod_{d\mid r}\Phi_d(X)
\]
remains a product of pairwise coprime polynomials because \(X^r-1\)
is separable. If the element were a zero of a factor indexed by a
proper divisor \(d\), its order would divide \(d\). Hence there
is a specialization \(R\to k\) sending the parameter to the
chosen \(s\).

Formula~\eqref{eq:source-integral} has no integer denominators.
To specialize \eqref{eq:invariance}, substitute the torus formulas
and clear only powers of
\[
x,\ y,\ t,\ s,\ sx-y.
\]
The resulting numerator is an element of an integral
polynomial/Laurent ring over \(R\). It is zero in the fraction
field by the source theorem, and is therefore zero in that ring.
Specialization preserves the identity of rational functions over
\(k\). The allowed phase parameters \(s,t\) are units, and
\(sx-y\) is not the zero polynomial after specialization. This
argument imports no smoothness assertion from characteristic zero.

\subsection{Extension across the accessible lines}

The function \(I_{r,t}\) is regular on the torus. The generic point
of \(L_4\) reaches the torus after one forward step. The generic
points of \(L_3,L_2,L_1\) do so after two, three and four steps,
respectively. To see this, use Table~\ref{tab:native}: each
intermediate map between line coordinates is affine with nonzero
slope, and the last \(L_4\) coordinate is generically nonzero.

All these forward maps are morphisms by
\eqref{eq:extension-torus}--\eqref{eq:extension-L4}. Repeated
invariance therefore expresses \(I_{r,t}\), near each line's
generic point, as the pullback of a regular Laurent function on a
target torus. It has no pole along any of the four accessible lines.
Every other prime divisor of \(U\) meets the torus, so it has no
pole there either. Smoothness makes \(U\) normal; on a normal
variety a rational function without a codimension-one pole is regular.
Consequently
\begin{equation}\label{eq:integral-regular}
I_{r,t}\in\Gamma(U_{t,s},\OO_{U_{t,s}}).
\end{equation}

This conclusion includes the special points with zero line coordinate.
There is no isolated pole of a rational function left at one of those
points. The two sides of \eqref{eq:invariance} are now regular
functions on \(U_{t,s}\) agreeing on a dense open subset, so they
agree everywhere on \(U_{t,s}\).

\subsection{Exact order on the boundary}

For completeness, reproduce the nonzero leading term calculated in
the proof of the source's Theorem~3.1 \citep{JoshiRoffelsen2026}.
At the generic point of \(x=0\), with \(y\) finite and nonzero,
\begin{equation}\label{eq:idempotent-leading}
A(z)=-\frac{tz}{x}V+O(1),\qquad
V=\begin{pmatrix}1&0\\1&0\end{pmatrix},\quad
V^2=V,\quad\Tr V=1.
\end{equation}
The coefficient of \(x^{-r}\) in the product defining
\eqref{eq:source-integral} is
\((-1)^rs^{r(r-1)/2}t^rV\). The product of all the roots
of \(X^r-1\) is \((-1)^{r-1}\); since the roots are
\(1,s,\ldots,s^{r-1}\), this gives
\[
(-1)^rs^{r(r-1)/2}=-1.
\]
Taking the trace yields
\begin{equation}\label{eq:nonzero-leading}
I_{r,t}=-t^rx^{-r}+O\bigl(x^{-(r-1)}\bigr).
\end{equation}
Its leading coefficient never vanishes on the allowed domain, even
in characteristic two. The integral is nonconstant and has exact
pole order \(r\) on \(D_5\), the strict transform of \(x=0\).

Apply Lemma~\ref{lem:uniform-poles} to
\eqref{eq:integral-regular}. The order on this one boundary component
determines the entire polar divisor:
\begin{equation}\label{eq:exact-fibration}
\polediv(I_{r,t})=rD,\qquad
f_t:S_{t,s}\longrightarrow\PP^1,\qquad
f_t^{-1}(\infty)=rD.
\end{equation}
The morphism is defined over \(k\), since the surface and function
are. Its nonconstancy and projectivity imply surjectivity onto
\(\PP^1\). Every finite fibre is a nonzero effective divisor
of class \(rD\), disjoint from \(D\).
